% \pagenumbering{arabic}
\Large
\label{Chapter_6}
\section{Conclusion}
\label{sec:conclusion}

This thesis proposed a trust-partitioning architecture that delivers trustworthy security for constrained, mission-critical industrial deployments by placing each security function on the layer whose cost profile sustains it rather than applying uniform protection across devices that cannot afford it. To establish the proposition, the study developed: (i) STRATA, a unified architecture separating high-frequency behavioural validation on a fast permissioned layer from low-frequency identity anchoring, audit, and revocation committed to a public layer; (ii) a lightweight federated behavioural-trust module built on a trimmed one-dimensional convolutional recurrent model coupled with a resource-optimised permissioned ledger; (iii) a quantum-resilient identity and credential module based on entropy-seeded derivation and compressed post-quantum key delivery; and (iv) a consolidated evaluation spanning detection quality, resource overhead, ledger cost, and quantum resilience across two distinct ledger substrates.

The architecture operates as follows. Behavioural evidence is produced at the edge by a model carrying 41{,}762 parameters at 0.17~MB, committed to the permissioned layer under authority-weighted consensus, aggregated by epoch, and anchored as a digest on a public substrate beyond the reach of the operating consortium. Device credentials are derived from entropy captured at the device and delivered under lattice-based encapsulation, and a trust score accumulated from behavioural evidence drives withdrawal of those credentials through two independent triggers, with the withdrawal published for relying parties who do not trust the operator.

The proposed architecture was evaluated against eight recent federated blockchain approaches and eight representative ledger-based data-sharing frameworks. The results demonstrate that the approach is both adequate and deployable, achieving detection accuracy between 96.84\% and 97.70\% with a false positive rate of 1.90\%, finalisation in 0.11~s without confirmation depth, ledger metadata occupying 0.13\% of the link the learning traffic already uses, a ciphertext reduction of 46.6\% under reuse-aware delivery, and revocation confirmed in 1.8~s with success above 95\% across a tenfold increase in device population. Detection accuracy is 2.27 percentage points below the strongest reported comparator, which attains that result with a model twenty-nine times larger and a memory footprint exceeding the budget of a microcontroller-class industrial node. The architecture is nonetheless the only entry in the comparison reporting ledger cost, contract assessment, and support for constrained nodes together, and the only framework providing device lifecycle management, intelligent threat detection, quantum-resilient cryptography, and public auditability simultaneously.

Each of the five research questions posed in Chapter~\ref{Chapter_1} is answered by this evidence, as Table~\ref{tab:rq-answers} records in full. Trust can be partitioned across a hybrid ledger within real device budgets (RQ1), since finalisation completes in 0.11~s while the ledger occupies 0.13\% of the uplink and public cost is governed by the rate of withdrawal rather than by traffic volume. Lightweight behavioural trust scoring operates inside the budget of edge-class hardware without sacrificing detection (RQ2), at 0.17~MB and 3.1~ms per inference with accuracy between 96.84\% and 97.70\%. Post-quantum identity and credential delivery is affordable on constrained links (RQ3), with the delivered ciphertext reduced from 768 to 410~bytes and the exchange shortened by 30.7\%. The partition holds on two substrate pairs that share no consensus, execution, or cost model (RQ4). Decay-based withdrawal isolates a misbehaving node within 1.8~s, succeeds for more than 95\% of detected anomalies across a tenfold increase in device population, and leaves detection quality unchanged (RQ5). Each answer holds within the scope stated beside it, and the measurements that would extend that scope are identified among the practical concerns that follow.

The experimental results further establish that trust-partitioning is a property of the design rather than of any particular ledger. Each architectural role was discharged on two substrate pairs sharing no consensus mechanism, no execution model, and no cost model, while the boundary carried digests and hashed identifiers alone in each case, and the cost relation derived at design time was exhibited by the measured public-chain prices on the realisation where commitments are priced. The security, scalability, and speed trilemma, which binds within a single distributed ledger, is accordingly dissolved across two, with each demand routed to the layer whose cost profile matches it.

\section{Practical Concerns and Future Works}
\label{sec:future-work}

Deploying a trust-partitioned security architecture on constrained industrial hardware raises several research and implementation concerns that warrant consideration for the wider viability of the approach.

\begin{enumerate}
\item \textbf{Entropy Sourcing on Fielded Hardware.} Identifier derivation consumes a sample from the device entropy source and treats it as an opaque bit string, so the construction is indifferent to which source supplies it. The present realisation is seeded from a software generator, and while the derivation, storage, and lifecycle operations behave as specified under that substitution, the unpredictability a fielded device would obtain remains to be characterised. Hardware sources differ substantially in availability across the industrial device population, from secure elements carrying a dedicated generator to legacy field devices offering only a sensor noise floor. Research establishing the entropy actually attainable on each class of device, together with conditioning procedures appropriate to the weaker sources, would place the unpredictability argument on measured rather than architectural ground.

\item \textbf{Cost of Signature Migration.} Key establishment has been migrated to lattice-based encapsulation, where the exposure is retroactive and migration cannot be deferred. Digital signatures carry a prospective exposure instead, since a signature cannot be forged retrospectively onto a commitment already published, which makes deferral defensible but not indefinite. Migration multiplies key and signature sizes at every point where a block is signed and every point where a credential is presented, and those artefacts are stored rather than transmitted once. Sizing that cost against the storage budget of a field gateway, and identifying which signature family fits it, is an engineering question the present design leaves open.

\item \textbf{Privacy of Model Updates Under Inference Attack.} The federated arrangement preserves data locality by ensuring that raw industrial traffic never leaves the device. Model updates are nonetheless exchanged without secure aggregation or noise injection, a choice made to preserve deployability on nodes whose budgets the evaluation showed to be tight. Within the evaluated threat model the exposure is bounded by authenticated participation, update validation, and ledger-recorded provenance, so inference would have to be attempted from an attributable position. Deployments operating under stronger adversarial privacy requirements would require secure aggregation or differential privacy, and establishing the accuracy and computation cost those mechanisms impose at the parameter counts involved is a direct extension of this work.

\item \textbf{Anchoring Economics Under Congestion.} Public anchoring cost separates into a term compressible by raising the epoch span and a term set by the rate at which devices are withdrawn. That separation was confirmed against measured prices under ordinary conditions, and it establishes that expenditure is governed by how often devices misbehave rather than by how much traffic they generate. Sustained congestion on the public substrate raises the per-commitment fee for both terms simultaneously, and the epoch span can absorb the first but not the second. Research into batching strategies for withdrawal publication, and into anchoring on cost-optimised layers, would extend the economic argument to conditions the present evaluation did not exercise.

\item \textbf{Characterisation of Admissible Substrates.} Realising the partition on two unrelated substrate pairs establishes that the design can be moved, and the properties claimed for it were obtained by machinery with nothing in common. Establishing that it can be moved to any substrate is a stronger claim requiring a characterisation of the substrate properties on which faithfulness depends, together with a proof that any substrate possessing them admits a faithful realisation. That direction also admits a negative formulation, namely whether classes of substrate exist on which no faithful realisation is possible, which would bound the generality of trust-partitioning from below as the two realisations bound it from above. Both formulations are open.

\item \textbf{Industrial Traffic Capture and Dataset Realism.} Evaluation drew on a public network benchmark, a general intrusion benchmark, and an internally generated industrial set constructed at the feature dimensionality the deployed model consumes, with agreement across the three establishing that detection does not depend on the peculiarities of any one. Synthetic generation reproduces the shape of industrial traffic and cannot reproduce the incidents, misconfigurations, and long-tailed device behaviour a live capture would contain. The architecture treats the detector as replaceable for this reason, since behavioural scoring is a stratum with a defined output and a deployment holding proprietary traffic substitutes its own model without disturbing the strata above it. Access to captured operational traffic, and evaluation of the partition under it, would close the remaining distance between benchmark performance and plant performance.

\item \textbf{Instrumented Measurement on Microcontroller-Class Hardware.} Inference latency and processor occupancy were measured on an ARM Cortex-A class environment, the post-quantum delivery figures were obtained on a workstation bench with single-board microbenchmarks, and energy per inference was estimated from an embedded-power model rather than read from an instrumented supply. The memory figures that decide feasibility follow from parameter count and encoding and would carry over unchanged, but latency, occupancy, and energy on a microcontroller with a lower clock and no floating-point unit would not. Repeating the behavioural and cryptographic measurements on instrumented microcontroller-class parts, and implementing the ledger bridge that binds revocation to the anchoring stratum in the identity realisation, would replace the remaining estimates with measurements at the lower end of the device spectrum the architecture targets.
\end{enumerate}

